\begin{table}[t]
    \centering
    \caption{Qwen2.5-0.5B inference performance. DQ denotes dynamic
    quantization with outlier correction. TTFT is the median of 10 warm
    requests with 128 input tokens, batch size 1, and four CPU threads;
    model loading and compilation are excluded.}
    \label{tab:architecture-performance}
    \footnotesize
    \setlength{\tabcolsep}{3pt}
    \begin{tabular}{@{}lrr@{}}
        \hline
        Execution configuration & TTFT (ms) & Peak memory (MiB) \\
        \hline
        Intel CPU, without DQ          & 293.17 & 2301.68 \\
        Intel CPU, with DQ             & 598.33 & 1306.64 \\
        AMD CPU, without DQ            & 331.40 & 2322.67 \\
        AMD CPU, with DQ               & 214.02 & 1383.15 \\
        AMD CPU+NPU, without DQ        & 725.41 & 1913.98 \\
        AMD CPU+NPU, with DQ           & 816.72 & 1609.84 \\
        Rocket+Gemmini, with DQ & TBD & TBD \\
        \hline
    \end{tabular}
    \vspace{2pt}
    \parbox{\columnwidth}{\scriptsize Peak memory is sampled inference-phase
    host RSS plus mapped NPU buffer allocation in shared DDR; the latter is
    excluded from RSS and is not a separate physical-residency measurement.
    CPU baselines use FP32. Single-core NPU matmuls use BF16 without DQ; DQ uses INT8 with
    FP32 CPU outlier correction. Rocket+Gemmini is not measured here.}
\end{table}
